\subsection{Observable Signatures}
  \label{subsec:void_signatures}

  The effective saturation framework leads to several testable
  observational signatures that distinguish it from both
  $\Lambda$CDM and purely local-void interpretations.

  First, the locally inferred Hubble constant should correlate
  with the large-scale density environment.
  Measurements performed in or near deep voids are expected to yield
  systematically higher values of $H_0$ than those obtained in
  denser regions.
  This correlation should persist after controlling for known
  observational systematics and distance-calibration effects.

  Second, the effect should exhibit a characteristic redshift
  dependence.
  At sufficiently large redshift, line-of-sight averaging over
  multiple environments suppresses the environmental bias.
  The inferred expansion rate should therefore converge toward
  $H_{\mathrm{global}}$, producing a gradual reduction of the
  offset with increasing survey depth.

  Third, the mechanism predicts distinct signatures in void-related
  observables, including weak lensing profiles and redshift drift
  measurements~\cite{Cai2017VoidLensing}.
  In particular, if the effective expansion rate inside voids is
  enhanced relative to the global value, the temporal evolution
  of redshift for sources embedded in deep underdensities should
  deviate from the homogeneous expectation in a systematic way.

  More generally, the predicted shift depends not only on the
  local environment of the observer, but on the integrated
  low-density structure sampled along the line of sight.
  Both the location of the source and the geometry of intervening
  voids contribute to the inferred bias.
  This implies that samples with different sky coverage or
  different environmental selection criteria may yield
  systematically different $H_0$ estimates.

  Importantly, the framework does not predict a unique numerical
  value for the offset in a given environment.
  Instead, once $a_\star$ is fixed independently from galactic
  rotation curves, the sign, scaling behavior, and environmental
  dependence of the effect are determined.
  This renders the proposal falsifiable with current and
  forthcoming large-scale-structure and distance-ladder data.
